\documentclass[12pt]{article}
\usepackage[a4paper, margin=1in]{geometry}
\usepackage{graphicx}
\usepackage{hyperref}
\usepackage{setspace}

\setcounter{secnumdepth}{0}

\title{\textbf{Agentic AI for Clinical Appointment No-Show Prediction}\\
\vspace{0.3cm}
\large Intelligent Care Coordination using RAG and LLM Agents}

\author{
\textbf{Team Members:} \\
Balanagu Sesha Srikar \\
Siraparapu Siva Sankar Mani Kanta \\
Shaik Junaid Sami
}

\date{\today}

\begin{document}

\maketitle
\onehalfspacing

%--------------------------------------------------

\section{Abstract}
This project presents an Agentic AI system designed for intelligent care coordination in healthcare operations. The system predicts patient appointment no-show risk and extends beyond prediction by incorporating reasoning, guideline retrieval, and structured intervention planning. Using a ReAct-style agent architecture with Retrieval-Augmented Generation (RAG), the system enables autonomous decision-making. An evaluation loop further improves reliability by validating outputs and enabling self-correction.

%--------------------------------------------------

\section{Introduction}
Missed clinical appointments significantly impact healthcare efficiency and patient outcomes. While predictive models estimate no-show probability, they fail to provide actionable insights.

This project introduces an Agentic AI system that not only predicts risk but also explains it, retrieves domain-specific guidelines, and generates structured intervention plans. The system demonstrates how modern AI can transition from prediction to decision support.

%--------------------------------------------------

\section{Problem Statement}
The objective is to build an intelligent system that:
\begin{itemize}
    \item Predicts appointment no-show risk
    \item Explains contributing risk factors
    \item Retrieves relevant healthcare guidelines
    \item Generates actionable intervention strategies
    \item Ensures ethical and reliable outputs
\end{itemize}

%--------------------------------------------------

\section{Agentic AI System}

\subsection{Agent Architecture}
The system uses a ReAct-style agent implemented using LangGraph. The agent dynamically selects tools to perform reasoning tasks instead of following a fixed pipeline.

\subsection{Tool Design}
The agent is equipped with domain-specific tools:
\begin{itemize}
    \item Prediction Tool (ML model)
    \item Risk Analysis Tool (LLM reasoning)
    \item Guideline Retrieval Tool (RAG)
    \item Intervention Planning Tool
    \item Risk Flag Detection Tool
\end{itemize}

\subsection{Retrieval-Augmented Generation (RAG)}
A structured knowledge base of healthcare guidelines is embedded into a vector database. The system retrieves relevant context using semantic similarity search, ensuring evidence-based recommendations.

\subsection{Evaluation Loop}
A secondary LLM evaluates outputs based on:
\begin{itemize}
    \item Clinical Accuracy
    \item Guideline Alignment
    \item Ethical Soundness
\end{itemize}

If the output fails evaluation, the system retries with feedback, enabling iterative refinement and improved reliability.

%--------------------------------------------------

\section{System Workflow}

The system follows an agent-driven workflow:

\begin{enumerate}
    \item Accept patient input data
    \item Predict no-show risk
    \item Analyze contributing factors
    \item Retrieve relevant guidelines (RAG)
    \item Generate structured intervention plan
    \item Evaluate output using critic model
    \item Retry if necessary
\end{enumerate}

%--------------------------------------------------

\section{Results}

The system produces:
\begin{itemize}
    \item Risk assessment summary
    \item AI-generated clinical reasoning
    \item Evidence-based guideline retrieval
    \item Structured intervention plan
\end{itemize}

The integration of RAG and evaluation loop improves output quality, reduces hallucinations, and ensures actionable recommendations.

%--------------------------------------------------

\section{Application Interface}

A Streamlit-based user interface is developed to:
\begin{itemize}
    \item Input patient details
    \item Display risk predictions
    \item Show reasoning and explanations
    \item Present retrieved guideline evidence
    \item Visualize intervention strategies
    \item Display evaluation results
\end{itemize}

%--------------------------------------------------

\section{Ethical Considerations}
The system ensures:
\begin{itemize}
    \item Transparency in decision-making
    \item Fair and unbiased recommendations
    \item Patient-centric interventions
\end{itemize}

Predictions are used as decision-support tools, not deterministic outcomes.

%--------------------------------------------------

\section{Future Scope}
\begin{itemize}
    \item Integration with real-time hospital systems
    \item Multi-modal AI capabilities
    \item Advanced personalization using patient history
    \item Scalable deployment across healthcare networks
\end{itemize}

%--------------------------------------------------

\section{Conclusion}
This project demonstrates the implementation of an Agentic AI system capable of reasoning, retrieval, planning, and evaluation. By combining machine learning with LLM-based intelligence, the system provides a comprehensive and practical solution for healthcare care coordination.

%--------------------------------------------------

\section{References}
\begin{itemize}
    \item LangChain Documentation
    \item LangGraph Documentation
    \item ChromaDB Documentation
    \item Groq API Documentation
\end{itemize}

\end{document}